\documentclass{article}

% ============================================================================
% Track selection
% ============================================================================
% Pick exactly ONE of the following lines. The default below = Main Track,
% double-blind reviewing. Switch by uncommenting a different `[opt]` line and
% commenting out the default `\usepackage{neurips_2026}`.
%
% \usepackage{neurips_2026}                          % main (default)
% \usepackage[main]{neurips_2026}                    % explicit Main Track
% \usepackage[position]{neurips_2026}                % Position Paper Track
% \usepackage[eandd]{neurips_2026}                   % Evaluations & Datasets
% \usepackage[creativeai]{neurips_2026}              % Creative AI Track
% \usepackage[sglblindworkshop]{neurips_2026}        % Workshop (single-blind)
% \usepackage[dblblindworkshop]{neurips_2026}        % Workshop (double-blind)
%
% For camera-ready, ADD `final` to the track option, e.g.:
%   \usepackage[main, final]{neurips_2026}
%
% For arXiv / preprint distribution:
%   \usepackage[preprint]{neurips_2026}
%
\usepackage{neurips_2026}

% For workshop submissions, also set:
% \workshoptitle{Your Workshop Name}

\usepackage[utf8]{inputenc}
\usepackage[T1]{fontenc}
\usepackage{hyperref}
\usepackage{url}
\usepackage{booktabs}
\usepackage{amsfonts}
\usepackage{nicefrac}
\usepackage{microtype}
\usepackage{xcolor}
\usepackage{graphicx}
\usepackage{amsmath}
\usepackage{amssymb}

% ============================================================================
% Title
% ============================================================================
\title{Your Paper Title Here}

% ============================================================================
% Authors
%
% By default, NeurIPS uses double-blind reviewing — author info is hidden
% from the compiled PDF until [final] is set. You can still fill the
% \author block; it only renders in the camera-ready.
% ============================================================================
\author{
  Alice Author \\
  Department of Computer Science \\
  Your University \\
  \texttt{alice@univ.edu}
  \And
  Bob Coauthor \\
  Department of Statistics \\
  Your University \\
  \texttt{bob@univ.edu}
}

\begin{document}
\maketitle

% ============================================================================
\begin{abstract}
% Abstract is ~200 words. State problem, gap, contribution, headline result.
% Avoid vague "we propose" — name what's specifically new.
Replace this paragraph with your abstract.
\end{abstract}

% ============================================================================
\section{Introduction}
\label{sec:intro}

Motivate the problem in one paragraph. State the gap in prior work in the
next. Then list contributions explicitly:

\begin{itemize}
    \item Contribution 1 --- be specific.
    \item Contribution 2.
    \item Contribution 3.
\end{itemize}

Close with a paragraph on paper organization (Section~\ref{sec:related}
reviews \dots; Section~\ref{sec:method} introduces \dots).

% ============================================================================
\section{Related Work}
\label{sec:related}

Group related work by theme, not by venue. For each closest paper, state
explicitly how this work differs.

% ============================================================================
\section{Method}
\label{sec:method}

State assumptions upfront. Define notation on first use.

% ============================================================================
\section{Experiments}
\label{sec:exp}

\subsection{Setup}
Datasets, baselines, metrics, hardware, hyperparameters, random seeds.

\subsection{Main Results}
\begin{table}[t]
\centering
\caption{Main results. Best in \textbf{bold}.}
\label{tab:main}
\begin{tabular}{lcc}
\toprule
Method & Metric A $\uparrow$ & Metric B $\downarrow$ \\
\midrule
Baseline 1 & 65.4 & 0.42 \\
Baseline 2 & 67.1 & 0.39 \\
\textbf{Ours} & \textbf{72.3} & \textbf{0.31} \\
\bottomrule
\end{tabular}
\end{table}

\subsection{Ablations}
One ablation per contribution claim in the introduction.

% Example figure:
% \begin{figure}[t]
% \centering
% \includegraphics[width=0.7\linewidth]{figures/teaser.pdf}
% \caption{One-line caption that stands alone.}
% \label{fig:teaser}
% \end{figure}

% ============================================================================
\section{Discussion}
\label{sec:discussion}

What did we learn? Where does the method break? Honest limitations belong
here — NeurIPS reviewers reward candor (and the Checklist explicitly asks
about limitations).

% ============================================================================
\section{Conclusion}
\label{sec:conclusion}

Restate the contribution (no inflation). Future work should be specific.

% ============================================================================
% Bibliography
% ============================================================================
% NeurIPS uses plainnat by default (loaded by the .sty). plainnat reads .bib
% with natbib-style citations.
\bibliographystyle{plainnat}
\bibliography{references}

% ============================================================================
% NeurIPS Paper Checklist — REQUIRED for submission.
% Papers without the checklist are desk-rejected.
%
% The checklist does NOT count toward the page limit. It goes AFTER the
% references and (optional) supplemental material.
% ============================================================================
\newpage
\input{checklist}

% ============================================================================
% Appendix (uncomment if needed; goes after checklist)
% ============================================================================
% \appendix
% \section{Additional Experiments}
% \label{app:additional}

\end{document}
